\subsection{瓶颈空间簇与容量影子价值模型}

任务负荷容量边界给出了系统最多能承受的任务规模，但仍需进一步定位限制调度性能的空间资源。

\subsubsection{边界工况与离散影子价值}

本节选择任务负荷与系统容量边界中接近容量边界且仍满足业务时域的工况，并在该工况上重新展开 P3 时空冲突修复。\textbf{我们把仓库局部通道视为有限容量资源，计算其容量单位放松后的边际价值。}

设 $C$ 为空间网格连通簇，$N_{C,t}$ 为时间桶 $t$ 中占用簇 $C$ 的 AMR 数。基线容量约束抽象为
\begin{equation}
N_{C,t}\leq q_C,\qquad q_C=1.
\end{equation}
这里的 \textbf{$q_C$ 是局部能同时通过的 AMR 数量}。\textbf{执行 $q_C:1\rightarrow 2$ 可以解释为通道拓宽、设置会车区、调整局部交通规则或增加分拣口前缓冲区。}由于当前 P2 的 MILP 没有显式建立边容量约束，本部分对候选空间簇逐一放松容量并重跑 P3 冲突修复优化，得到\textbf{离散影子价值}
\begin{equation}
\pi_C^{J}=J^{\mathrm{base}}-J^{+1}(C),\qquad
\pi_C^{T}=\Cmax^{\mathrm{base}}-\Cmax^{+1}(C),
\end{equation}
其中 P3 修复目标为
\begin{equation}
J=100D_{\mathrm{sum}}+5\Cmax+W+2R+20S+20W_0.
\end{equation}
$D_{\mathrm{sum}}$ 为总延期，$W$ 为冲突修复插入等待，$R$ 为换路次数，$S$ 为同车换序次数，$W_0$ 为初始等待。$\pi_C^J$ 和 $\pi_C^T$ 分别表示簇容量增加 1 后目标函数和最大完工时间的真实改善量；\textbf{最终瓶颈排序以重优化后的有限差分为准，而不是以经过次数或热图颜色为准。}

\begin{center}
\small
\resizebox{0.98\linewidth}{!}{%
\begin{tikzpicture}[
    font=\small,
    analysisstep/.style={
        draw=blue!28!black,
        fill=blue!4,
        rounded corners=2pt,
        line width=0.55pt,
        minimum height=1.18cm,
        text width=3.25cm,
        align=center,
        inner sep=5pt,
        text=black!88
    },
    decisionsstep/.style={
        draw=teal!24!black,
        fill=teal!3,
        rounded corners=2pt,
        line width=0.55pt,
        minimum height=1.18cm,
        text width=3.25cm,
        align=center,
        inner sep=5pt,
        text=black!88
    },
    arrow/.style={-{Stealth[length=2.15mm]}, line width=0.55pt, draw=black!45},
    node distance=0.72cm and 0.92cm
]
\node[analysisstep] (base) {\textbf{1 基线工况}\\边界任务规模\\记录 $\Cmax^{\mathrm{base}},J^{\mathrm{base}},W^{\mathrm{base}}$};
\node[analysisstep, right=of base] (score) {\textbf{2 时空投影}\\轨迹映射到网格--时间桶\\计算 $s(c)$};
\node[analysisstep, right=of score] (cluster) {\textbf{3 候选簇}\\连通合并\\保留冲突/等待信号};
\node[decisionsstep, below=of cluster] (relax) {\textbf{4 容量放松}\\逐簇执行 $q_C\leftarrow q_C+1$\\重跑 P3 修复};
\node[decisionsstep, left=of relax] (value) {\textbf{5 影子价值}\\计算 $\pi_C^J,\pi_C^T$\\按 $\Delta\Cmax$ 排序};
\node[decisionsstep, left=of value] (decision) {\textbf{6 工程判读}\\识别主瓶颈\\给出布局/规则改造};
\draw[arrow] (base) -- (score);
\draw[arrow] (score) -- (cluster);
\draw[arrow] (cluster) -- (relax);
\draw[arrow] (relax) -- (value);
\draw[arrow] (value) -- (decision);
\end{tikzpicture}}
\captionof{figure}{瓶颈空间簇容量灵敏度分析流程}
\label{fig:stage4-sensitivity-flow}
\end{center}

\subsubsection{前置筛选方法}

我们对候选簇进行空间统计。轨迹点按 $1\mathrm{m}\times 1\mathrm{m}$ 网格和 $1\mathrm{s}$ 时间桶投影，并对同一 AMR 在同一时间桶内的重复采样去重。对每个网格 $c$ 计算占用、冲突修复事件、冲突等待和延期压力的标准化分量：
\begin{equation}
s(c)=0.25z(\mathrm{occ}_c)+0.30z(\mathrm{conf}_c)+0.30z(\mathrm{wait}_c)+0.15z(\mathrm{delay}_c).
\end{equation}
相邻高分网格合并成连通簇后，只对含有冲突或等待信号的簇做容量放松重求。\textbf{若一个区域只是经过次数高，但容量放松后 $\pi_C^J$ 近似为 0，则只能称为高流量区，不能作为主瓶颈。}
